\section{Results}

\subsection{Structural change}

Table~\ref{tab:structural} and Figure~\ref{fig:instrcounts} show the IR
instruction counts for the baseline, the framework, and \texttt{opt -O1}.
Figure~\ref{fig:reduction} shows the reduction the framework achieves relative to
the SSA baseline. The framework removes the most on \texttt{arith\_redundant},
at \ArithRedundantReduction\%, where the planted redundancy and identities are
exactly what it targets. On \texttt{clean} it makes no change at all, a
\CleanReduction\% reduction, which is the correct and expected result for a
kernel with nothing to simplify.

\begin{table}[h]
  \centering
  \begin{tabular}{lrrrr}
\toprule
kernel & baseline & my-default & reduction & opt -O1 \\
\midrule
arith\_redundant & 39 & 32 & 17.9\% & 34 \\
branchy & 42 & 39 & 7.1\% & 44 \\
clean & 27 & 27 & 0.0\% & 43 \\
dead\_code & 29 & 25 & 13.8\% & 14 \\
\bottomrule
\end{tabular}

  \caption{IR instruction counts by kernel and configuration, with the
    framework's reduction relative to the baseline.}
  \label{tab:structural}
\end{table}

\begin{figure}[h]
  \centering
  \includegraphics[width=0.85\textwidth]{instruction_counts.pdf}
  \caption{Instruction count by kernel and configuration.}
  \label{fig:instrcounts}
\end{figure}

\begin{figure}[h]
  \centering
  \includegraphics[width=0.85\textwidth]{instruction_reduction.pdf}
  \caption{Instruction reduction achieved by the framework, relative to the SSA
    baseline.}
  \label{fig:reduction}
\end{figure}

Figure~\ref{fig:opcode} breaks the change down by opcode for the showcase kernel,
making visible that the reduction comes from removing multiplies (strength
reduced to shifts and then deduplicated) and the redundant arithmetic.

\begin{figure}[h]
  \centering
  \includegraphics[width=0.85\textwidth]{opcode_delta.pdf}
  \caption{Per opcode counts for \texttt{arith\_redundant}, baseline against the
    framework.}
  \label{fig:opcode}
\end{figure}

\subsection{Runtime}

Table~\ref{tab:runtime} and Figure~\ref{fig:runtime} report median runtime with
the interquartile range as an error bar. On \texttt{arith\_redundant} the
framework moves the median from \ArithRedundantBaseMs{} ms to
\ArithRedundantOursMs{} ms, a difference well outside the spread and therefore
real signal. On \texttt{branchy} and \texttt{clean} the runtime change is within
the interquartile range and is reported as structural only, with the noise stated
rather than hidden.

\begin{table}[h]
  \centering
  \begin{tabular}{llrr}
\toprule
kernel & configuration & median (ms) & IQR (ms) \\
\midrule
arith\_redundant & baseline & 64.28 & 4.66 \\
arith\_redundant & my-default & 47.44 & 1.46 \\
arith\_redundant & opt -O1 & 29.45 & 1.28 \\
\addlinespace
branchy & baseline & 57.66 & 0.91 \\
branchy & my-default & 55.74 & 1.31 \\
branchy & opt -O1 & 23.24 & 0.17 \\
\addlinespace
clean & baseline & 62.20 & 5.91 \\
clean & my-default & 62.32 & 1.64 \\
clean & opt -O1 & 50.44 & 2.39 \\
\addlinespace
dead\_code & baseline & 42.37 & 6.18 \\
dead\_code & my-default & 43.12 & 7.44 \\
dead\_code & opt -O1 & 1.19 & 0.18 \\
\addlinespace
\bottomrule
\end{tabular}

  \caption{Median runtime and interquartile range by kernel and configuration.}
  \label{tab:runtime}
\end{table}

\begin{figure}[h]
  \centering
  \includegraphics[width=0.85\textwidth]{runtime_medians.pdf}
  \caption{Median runtime by kernel and configuration, with the interquartile
    range as an error bar. The backend was invoked at \texttt{-O0} so that middle
    end differences reach the binary.}
  \label{fig:runtime}
\end{figure}

\subsection{Correctness}

Table~\ref{tab:checksum} records that on every kernel the baseline, framework,
and \texttt{opt -O1} binaries print an identical checksum, so no transformation
changed observable behavior.

\begin{table}[h]
  \centering
  \begin{tabular}{lll}
\toprule
kernel & checksum & configurations agree \\
\midrule
arith\_redundant & \texttt{51247235914731008} & yes \\
branchy & \texttt{3361102545920} & yes \\
clean & \texttt{17182339206778112} & yes \\
dead\_code & \texttt{1331200000000} & yes \\
\bottomrule
\end{tabular}

  \caption{Differential execution: every configuration produces the same
    checksum per kernel.}
  \label{tab:checksum}
\end{table}
